%!TeX root=MemoriaTFG.tex

\chapter{Declaración sobre el uso de herramientas de inteligencia artificial}
\label{anexo:declaracion-ia}
\label{cap:anexod}

Este anexo documenta el uso de herramientas de inteligencia
artificial durante la elaboración del presente Trabajo de Fin de
Grado, en cumplimiento de las directrices institucionales de la
Universitat de les Illes Balears sobre transparencia académica en
el uso de tecnologías generativas.

\section{Alcance de la declaración}
\label{sec:anexod-alcance}

La declaración cubre la totalidad del periodo de redacción y
desarrollo del presente trabajo (curso académico 2025-2026). Se
distinguen tres ámbitos de uso, cuya descripción detallada se
desarrolla en las secciones siguientes: uso académico (redacción
del documento), uso técnico (desarrollo del código y de los
experimentos) y uso documental (organización de notas y
trazabilidad).

\section{Herramientas empleadas}
\label{sec:anexod-herramientas}

La tabla~\ref{tab:anexod-herramientas} enumera las herramientas
de inteligencia artificial empleadas durante el desarrollo del
trabajo, su modalidad de uso y el ámbito al que se aplican.

\begin{table}[H]
\centering
\footnotesize
\setlength{\tabcolsep}{4pt}
\renewcommand{\arraystretch}{0.95}
\caption{Herramientas de IA empleadas durante la elaboración del
TFG.}
\label{tab:anexod-herramientas}
\begin{tabular}{llll}
\hline
Herramienta              & Proveedor   & Modalidad           & Ámbito \\
\hline
Claude (Sonnet, Opus)    & Anthropic   & Asistente conversacional & Redacción, código, organización \\
GPT-4o, GPT-4o-mini      & OpenAI      & API programática      & Evaluación experimental (objeto del trabajo) \\
GPT-5 (familia)          & OpenAI      & API programática      & Razonamiento clínico del prototipo \\
Gemini~2.5 Pro/Flash     & Google      & API programática      & Evaluación experimental (objeto del trabajo) \\
MedGemma~4B y~27B        & Google      & Modelos locales       & Evaluación + razonamiento del prototipo \\
GitHub Copilot           & GitHub      & Asistencia en IDE     & Sugerencias de código \\
\hline
\end{tabular}
\end{table}

\section{Uso académico: redacción del documento}
\label{sec:anexod-academico}

Las herramientas conversacionales (Claude Sonnet y Opus) se han
utilizado durante la redacción del presente documento en cuatro
modalidades: (i)~revisión estilística y de coherencia narrativa
sobre el texto producido por el autor;
(ii)~reformulación de párrafos con tono no académico hacia un
tono paper estricto, conforme a las indicaciones explícitas del
director del trabajo;
(iii)~verificación cruzada entre afirmaciones cuantitativas del
texto y los reportes técnicos en el repositorio (auditoría
descrita en secci\'on~\ref{sec:anexod-verificacion});
(iv)~asistencia en la generación de tablas LaTeX a partir de
datos preexistentes en formato Markdown.

En ningún caso se ha utilizado la herramienta para generar
contenido original sobre los resultados experimentales, la
interpretación de los hallazgos o la formulación de las
conclusiones. Las cifras reportadas en el
capítulo~\ref{cap:resultados} y discutidas en el
capítulo~\ref{cap:discusion} provienen exclusivamente de los
experimentos ejecutados por el autor, documentados en los
reportes técnicos del repositorio
(secci\'on~\ref{sec:anexoa-tareas}) y verificables de forma
independiente.

\section{Uso técnico: desarrollo del código}
\label{sec:anexod-tecnico}

Las herramientas de asistencia en IDE (GitHub Copilot) se han
utilizado durante el desarrollo del código en dos modalidades:
(i)~autocompletado de sentencias rutinarias (carga de datasets,
formateo de salidas, configuración de \emph{loggers});
(ii)~generación inicial de plantillas de \emph{tests} unitarios y
de \emph{wrappers} de modelos a partir de la documentación
oficial de las bibliotecas empleadas.

El código de los protocolos experimentales centrales
(sondeo lineal, ajuste fino, segmentación, evaluación zero-shot,
auditoría de leakage) ha sido desarrollado por el autor con
referencia explícita a la implementación original publicada por
los autores de PanDerm~\cite{panderm2025} y de
DermLIP~\cite{derm1m}. La revisión de la corrección de los
algoritmos y la validación de los resultados se ha realizado por
el autor mediante reproducción cruzada de cifras de referencia
(p.~ej.~ajustar los hiperparámetros hasta reproducir el rango de
AUROC reportado por los autores sobre los datasets compartidos).

\section{Uso documental: organización y trazabilidad}
\label{sec:anexod-documental}

Las herramientas conversacionales se han utilizado adicionalmente
para tareas de organización documental que no afectan al contenido
del trabajo: redacción de notas internas, estructuración del
índice del repositorio, formateo de documentación auxiliar
(no incluida en el documento final) y generación de informes de
estado del proyecto durante el desarrollo.

\section{Modelos como objeto de evaluación}
\label{sec:anexod-objeto}

GPT-4o, GPT-4o-mini, Gemini~2.5~Pro, Gemini~2.5~Flash,
MedGemma~4B Vision y MedGemma~27B Text constituyen
\emph{objetos de evaluación} del presente trabajo
(capítulo~\ref{cap:resultados}, secci\'on~\ref{sec:llm-results}), no
herramientas auxiliares de redacción. Su uso se documenta en los
protocolos experimentales correspondientes
(secci\'on~\ref{sec:llm}). La declaración del presente anexo se refiere
exclusivamente al uso de inteligencia artificial como
\emph{instrumento de apoyo} a la elaboración del documento, no
como objeto experimental.

\section{Verificación cruzada de cifras y citas}
\label{sec:anexod-verificacion}

Durante la fase final de redacción se ha aplicado un protocolo
sistemático de verificación cruzada entre las afirmaciones
cuantitativas del texto y los reportes técnicos del repositorio.
El procedimiento, descrito en el documento interno
\texttt{PRECAUCIONES\_ANEXOS.md}, consiste en: (i)~identificación
de cada cifra cuantitativa del documento;
(ii)~localización de la fuente verificable
(\texttt{RESULTADOS\_TFG.md}, archivos en \texttt{results/} o
\texttt{output/}); (iii)~clasificación del estado en cuatro
categorías: \textsc{ok} (cifra verificada), \textsc{parcial}
(existe el mecanismo pero falta la ejecución reportada),
\textsc{falta} (sin evidencia) y \textsc{contradicción}
(divergencia entre texto y fuente). Las inconsistencias detectadas
se han corregido antes del depósito o documentado explícitamente
como \textsc{[no especificado]} cuando la cifra exacta requería
ejecución pendiente.

\section{Responsabilidad sobre el contenido}
\label{sec:anexod-responsabilidad}

La autoría intelectual del trabajo, la formulación de la pregunta
de investigación, el diseño experimental, la ejecución de los
experimentos, la interpretación de los resultados, la redacción
final del documento y la responsabilidad sobre las afirmaciones
contenidas corresponden íntegramente al autor del presente
Trabajo de Fin de Grado, bajo la supervisión académica del
director Dr.~Javier Varona Gómez y con la colaboración clínica de
la Dra.~Rosa Taberner Ferrer. Las herramientas de inteligencia
artificial empleadas constituyen instrumentos de apoyo cuya
contribución se enmarca en los términos descritos en este anexo.
